%==============================  program skeleton  =============================
\documentclass[11pt]{article}
\usepackage{subfiles}

\usepackage{geometry}
\usepackage[parfill]{parskip}
\usepackage{mathtools}
\usepackage{tcolorbox}
\usepackage{amsmath,amssymb,amsthm}              
\usepackage{graphicx}
\usepackage{enumitem}
\usepackage{tikz}
\usepackage{mathrsfs}  
\usetikzlibrary{calc}
\usepackage{tikz-cd}
\usepackage{xcolor}
\usepackage{booktabs}
\usepackage[linesnumbered,ruled,vlined]{algorithm2e} 
\usepackage{lineno}
\usepackage[colorlinks]{hyperref}
\usepackage{mdframed}
\usepackage[symbols,nogroupskip,sort=none]{glossaries-extra}
\renewcommand*{\glspostdescription}{\medskip}

\newtheorem*{theorem*}{Theorem}
\newtheorem*{proposition*}{Proposition}
\newtheorem*{corollary*}{Corollary}

\newcommand{\N}{\mathbb{N}}
\newcommand{\Z}{\mathbb{Z}}
\newcommand{\Q}{\mathbb{Q}}
\newcommand{\R}{\mathbb{R}}
\newcommand{\C}{\mathbb{C}}
\newcommand{\G}{\mathbb{G}}
\newcommand{\delayX}{\mathcal{X}}
\newcommand{\X}{\text{x}}
\newcommand{\W}{\text{w}}
\newcommand{\Rho}{\mathrm{P}}
\newcommand{\Lgr}{\mathcal{L}}
\mathchardef\Re="023C
\newcommand{\GLdelay}{\mathrm{GL}_{\mathrm{delay}}}    
\newcommand{\Hrose}{\mathscr{H}} 
\newcommand{\krose}{k_{\theta,\sigma}}


\theoremstyle{plain} 
\newtheorem{theorem}{Theorem}[section]
\newtheorem{lemma}{Lemma}[section]
\newtheorem{corollary}{Corollary}[section]
\newtheorem{proposition}{Proposition}[section]


\theoremstyle{definition} 
\newtheorem{definition}{Definition}[section]

\theoremstyle{remark} 
\newtheorem{remark}{Remark}[section]
\newtheorem{example}{Example}[section]
\newtheorem{assumption}{Assumption}[section]

\DeclareMathOperator{\diag}{diag}
\newcommand{\Var}{\mathrm{Var}}
\newcommand{\Cov}{\mathrm{Cov}}
\DeclareMathOperator{\E}{\mathbb{E}}

% a grey framed box with small padding
\newmdenv[
  linecolor=gray!60,
  linewidth=.8pt,
  topline=true,
  bottomline=true,
  skipabove=\baselineskip,
  skipbelow=\baselineskip,
  innerleftmargin=6pt,
  innerrightmargin=6pt,
  innertopmargin=4pt,
  innerbottommargin=4pt,
  backgroundcolor=gray!5,
  frametitlefont=\normalfont\bfseries,
]{assumptionbox}


% ------------------------------------------------------------------
% A numbered list tailored for assumptions:
%   (W1) ..., (W2) ...   or   (K1) ..., (K2) ...
% ------------------------------------------------------------------
\newlist{assumptionlist}{enumerate}{1}          % base = enumerate, depth = 1
\setlist[assumptionlist]{
    leftmargin=1.4cm,                          % indent
    labelsep=0.2cm,                            % space after label
    itemsep=0.4\baselineskip,                  % vertical gap between items
    label=\textbf{(\Alph*)}                    % default label (over-ridden per list)
}

\pagecolor{white}   % put your \newcommand definitions here

\title{The Rose Surface: Discrete Onsager--Machlup, Dial Graphs, and a Program for Data‑Driven Statistical Dynamics}
\author{Patrick S.\ Mahon}
\date{\today}
\begin{document}
\maketitle
\tableofcontents

%------------------------------------------------------------------ 
\section*{Guide to the two–layer release}
\begin{enumerate}[label=\textbf{\Alph*)},leftmargin=*]
\item \textbf{Main Letter (12–15 pp)}  
      Appears separately as \emph{Geometry, Entropy, and Operator Learning
      with Tunable Kernels}.  It contains the core bridge
      theorem, KL = least‑action, and one synthetic experiment.

\item \textbf{This Program Paper (full version)}  
      Provides historical survey §1, proofs, extended risk analysis,
      dial‑graph curvature, and open problems (Sections 2–7 below).
\end{enumerate}
\vspace{1ex}
\noindent
\emph{To avoid duplication:} we load the Letter as an external file,
then append the expanded material.

\subfile{../letter/body.tex}   % ---  your concise 12‑page letter  ---

%==================================================================
\section{Historical Survey of Discrete Onsager–Machlup}
\phantomsection        % keeps hyperref anchors even while empty
\addcontentsline{toc}{subsection}{TODO write narrative}
% TODO: 3–4 pages summarising Schnakenberg cycles, Kaiser–Jack–Zimmer,
%       GENERIC, etc.  Use Table 1 from literature PDF.

%------------------------------------------------------------------
\section{Proofs for Main Theorems}
\subsection{Proof of Proposition~\ref{prop:CGF}}
\phantomsection
% TODO: full Taylor–Itô calculation, cite Lemma B.1

\subsection{Proof of Theorem~\ref{thm:bridge}}
\phantomsection
% TODO: detailed carré‑du‑champ expansion, Kuwada identity derivation

\subsection{Proof of Theorem~\ref{thm:KL_eq_action}}
\phantomsection
% TODO: edgewise gradient ⇒ global gradient ⇒ KL = OM

%------------------------------------------------------------------
\section{Finite‑Sample Theory}
\subsection{Goldenshluger–Lepski Oracle Inequality}
\phantomsection
% TODO: disclose constant, variance bound lemma, covering argument

\subsection{Synthetic and Real‑Data Experiments}
\phantomsection
% TODO: double‑well Langevin, Lorenz‑63, delay embedding plots

%------------------------------------------------------------------
\section{Dial‑Graph Geometry}
\subsection{Directed Discrete Curvature}
\phantomsection
% TODO: Ollivier vs Bakry–Émery comparison; include Figure 5

\subsection{Helmholtz Decomposition and Fisher Clock}
\phantomsection
% TODO: proof of Lemma~\ref{lem:graph-Helmholtz}, numeric toy example

%------------------------------------------------------------------
\section{Programmatic Directions}
\begin{enumerate}[label=\textbf{P\arabic*},leftmargin=*]
\item **Intrinsic Fisher Clocks** in out‑of‑equilibrium networks.  
\item **Multiparameter Persistent Homology** for dial graphs.  
\item **Generalised GENERIC Operators** on rough data manifolds.  
\end{enumerate}

%------------------------------------------------------------------
\section{Concluding Outlook}
\phantomsection
% TODO: 1‑page roadmap tying together P1–P3.

%=======================  bibliography  ===========================
\bibliographystyle{plainnat}
\bibliography{rose_refs}
\end{document}
%==============================================================================
